\chapter{组合与离散结构：数数、连线和算法}

\section{一场联欢会前的点票难题}

初二（3）班要办联欢会，班长小林负责两件事。第一件：全班 25 名同学排一列队伍上台合影，班主任想知道“一共有多少种不同的排法”，好决定要不要在方案里写上“随机排队”。第二件：联欢会结尾有个传统节目，每位同学都要和其他每位同学握一次手，不能重复。主持人要控制时间，得先知道一共要握多少次手。

小林觉得第一件事不难：25 个人排队，大概是“几十种”“几百种”吧？第二件事他也想估算：每个人握 24 次手，25 个人就是 $25\times 24=600$ 次？

这两个问题看起来微不足道，却藏着本章的全部入口。第一个问题的答案会让你吓一跳：它不是几十种，也不是几百种，而是一个 26 位数。第二个问题的直觉算法$25\times24$犯了一个隐蔽的错误——把每一次握手数了两遍。你会发现，“数数”这件我们从幼儿园就开始做的事，一旦规模稍大、结构稍复杂，就会变成一个需要方法和定理的数学分支。这个分支叫\emph{组合数学}，它和它的近亲\emph{图论}、\emph{算法}一起，构成所谓“离散数学”——研究一个一个、分得清清楚楚的对象的数学。

为什么单独给“离散”立一个分支？这本书前面大部分的舞台是连续的：实数轴上的点、光滑的曲线、连续变化的速度。可生活里有太多东西不是连续的：同学是一个一个的，握手是一次一次的，道路网是一个路口接一个路口的，密码是一位一位的，计算机里的数据更是一条一条的 0 和 1。对付这些“数得清个数”的对象，需要另一套工具。本章的任务就是把这套工具交到读者手里：怎样不漏不重地数数，怎样把关系画成点和线，以及怎样让一台只会按步骤办事的机器替我们高效地走完这些点和线。

\section{先猜一猜，再错一次}

先回到握手问题。小林的算法是：每人握 24 次，25 人共 $25\times 24=600$ 次。请读者先凭直觉判断：这个数偏大还是偏小？

再用一个更小的班级检验。假设全班只有 3 个人：甲、乙、丙。按小林的算法，握手次数是 $3\times 2=6$ 次。可实际数一数：甲乙、甲丙、乙丙，只有 3 次。6 恰好是 3 的两倍。再看 4 个人的班级：甲乙、甲丙、甲丁、乙丙、乙丁、丙丁，共 6 次，而小林的算法给出 $4\times 3=12$，又是两倍。

问题出在哪？“甲和乙握手”与“乙和甲握手”是同一件事，但在“每人握 24 次”的算法里，这件事在甲的账上记了一次，又在乙的账上记了一次。每一次握手都有两只手，于是总账翻了一倍。把 600 除以 2，得到 300，这才是正确答案。

\begin{fanli}
\textbf{看似合理的说法}：$n$ 个人两两握手，每人握 $n-1$ 次，所以总次数是 $n(n-1)$。

\textbf{错在哪里}：每一次握手属于两个人，“逐人统计再相加”必然把每次握手数两遍。凡是“数一件事有多少种”的问题，先要问：我数出来的每一件，会不会其实是同一件被重复登记了？辨别“有序”与“无序”、分辨“重复计数”，是组合计数里最常见的翻车现场。同样的错误还有：20 支球队单循环比赛（每两队踢一场），场次不是 $20\times 19=380$，而是 $380\div 2=190$。
\end{fanli}

再看排队问题。25 人排成一列，有多少种排法？如果读者猜“几百种”，那就掉进另一个陷阱：低估了乘法的力量。我们一步步看。第一个位置有 25 种选择；选好之后，第二个位置剩 24 种选择；接着 23 种、22 种……直到最后一个位置只剩 1 种。每个位置的选择互不干扰地串联起来，总数是
\[
25\times 24\times 23\times\cdots\times 2\times 1.
\]
这个连乘积大到需要专门造一个记号：记作 $25!$，读作“25 的阶乘”。它的值是 $15511210043330985984000000$，一个 26 位数。如果每秒钟试一种排法，试完所有排法需要的时间，比宇宙的年龄还要长几千亿倍。阶乘的增长速度，是本章后面讨论“算法到底跑不跑得动”的关键角色。

\section{数数的两条军规与两件武器}

上面的两次失败提示我们：数数不能靠蛮力列清单，要靠结构。组合数学的军规只有两条。

\textbf{加法原理}：做一件事有若干类互不相交的办法，第一类有 $a$ 种，第二类有 $b$ 种，那么总共有 $a+b$ 种。\textbf{乘法原理}：做一件事要分几步完成，第一步有 $a$ 种做法，第二步有 $b$ 种做法，那么总共有 $a\times b$ 种。排队用乘法，因为 25 个位置是“串联”的；而“从家到学校可以坐 3 路公交中的一路，或者骑 2 条路线中的一条”用加法，因为公交和骑车是“并联”的类别。

军规之下，最常用的两件武器是\emph{排列}和\emph{组合}。

从 $n$ 个人里挑出 $k$ 个人排成一列（挑出来还要分先后次序），排法数记作
\[
A_n^k=n(n-1)(n-2)\cdots(n-k+1)=\frac{n!}{(n-k)!}.
\]
例如从 8 名同学里选出场次表的冠、亚、季军，有 $A_8^3=8\times 7\times 6=336$ 种结果。如果不要名次、只要“挑出 3 个人组成组委会”，那么每 3 个人的小组会在排列里被重复计算 $3!=6$ 次（3 个人内部有 6 种排队方式），所以要除以 6。这样得到的数记作
\[
\binom{n}{k}=\frac{n!}{k!\,(n-k)!},
\]
读作“$n$ 选 $k$”，也叫\emph{组合数}或\emph{二项式系数}。从 8 人里选 3 人组委会，有 $\binom{8}{3}=56$ 种。握手问题用它一句话解决：25 人两两握手，次数就是 $\binom{25}{2}=\dfrac{25\times 24}{2}=300$。

\begin{shuyu}{组合数}
\textbf{直观解释}：从 $n$ 个东西里不管顺序地抓出 $k$ 个，有多少种抓法。

\textbf{正式定义}：$\dbinom{n}{k}=\dfrac{n!}{k!(n-k)!}$，其中 $0\le k\le n$，约定 $0!=1$。

\textbf{一个例子}：$\dbinom{5}{2}=\dfrac{5\times4}{2}=10$，即 5 个篮球队两两比赛共 10 场。
\end{shuyu}

组合数有一张著名的“地图”。把它们按行排开，每一行两端的数都是 1，中间的每个数等于它肩上两个数之和：
\[
\begin{array}{ccccccc}
& & & 1 & & &\\
& & 1 & & 1 & &\\
& 1 & & 2 & & 1 &\\
1 & & 3 & & 3 & & 1\\
& 1 & & 4 & & 6 & \quad\cdots
\end{array}
\]
第 4 行（从第 0 行数起）是 $1,4,6,4,1$，正好是 $\binom{4}{0},\binom{4}{1},\binom{4}{2},\binom{4}{3},\binom{4}{4}$。这张表在中国叫杨辉三角，在西方叫帕斯卡三角，数学家们各自独立地反复发现过它。它的构造规则本身就是一个计数论证：从 $n$ 个东西里选 $k$ 个，要么选中了某个特定的东西（剩下 $k-1$ 个从 $n-1$ 个里选），要么没选中（$k$ 个全从 $n-1$ 个里选），所以
\[
\binom{n}{k}=\binom{n-1}{k-1}+\binom{n-1}{k}.
\]
这个等式有一个值得注意的味道：它用一个更小规模的同类型问题来回答当前问题。这种“自己调用自己的缩小版”的想法叫\emph{递推}，是本章的第三件武器。

递推最经典的例子是爬楼梯：一段楼梯有 $n$ 级，每步可以跨 1 级或 2 级，有多少种走法？设走法数为 $f_n$。最后一步只有两种可能：从第 $n-1$ 级跨 1 级上来，或者从第 $n-2$ 级跨 2 级上来，所以
\[
f_n=f_{n-1}+f_{n-2},\qquad f_1=1,\ f_2=2.
\]
于是 $f_3=3$，$f_4=5$，$f_5=8$，$f_6=13$……这就是著名的斐波那契数列。注意它和杨辉三角的相似之处：一个看似要“枚举所有走法”的吓人问题，被一句“最后一步从哪里来”拆成了两个更小的问题。这种拆法后面会在“递归算法”一节里再次登场。

为什么组合数又叫“二项式系数”？把 $(a+b)^n$ 展开试试。$(a+b)^2=a^2+2ab+b^2$，系数是 $1,2,1$；$(a+b)^3=a^3+3a^2b+3ab^2+b^3$，系数是 $1,3,3,1$——正是杨辉三角的行。这不是巧合：展开 $(a+b)^n$ 时，$a^k b^{\,n-k}$ 这一项，相当于在 $n$ 个因子里挑出 $k$ 个贡献 $a$、其余贡献 $b$，挑法恰好是 $\binom{n}{k}$ 种。于是我们得到一个核心结论。

\begin{dingli}[二项式定理]
对任意数 $a,b$ 和非负整数 $n$，有
\[
(a+b)^n=\binom{n}{0}a^n+\binom{n}{1}a^{n-1}b+\binom{n}{2}a^{n-2}b^2+\cdots+\binom{n}{n}b^n.
\]
\end{dingli}

\begin{proof}
把 $(a+b)^n$ 写成 $n$ 个因子 $(a+b)$ 的连乘。展开时，每个因子要么贡献一个 $a$，要么贡献一个 $b$。要得到 $a^k b^{\,n-k}$，必须在 $n$ 个因子中恰好选出 $k$ 个来贡献 $a$，这样的选法有 $\binom{n}{k}$ 种，所以该项的系数是 $\binom{n}{k}$。把 $k=0,1,\dots,n$ 的所有项加起来，就得到公式。
\end{proof}

令 $a=b=1$，立刻得到一个数值例子：$\binom{4}{0}+\binom{4}{1}+\binom{4}{2}+\binom{4}{3}+\binom{4}{4}=2^4=16$。也就是说，一个 4 人小组的所有可能的“子集代表团”（从 0 人代表团到 4 人全去）一共 16 种。读者可以直接验证：$\binom{4}{0}=1$、$\binom{4}{1}=4$、$\binom{4}{2}=6$、$\binom{4}{3}=4$、$\binom{4}{4}=1$，加起来确实是 16。

最后一件数数武器简单得不像定理，却威力惊人。

\begin{dingli}[鸽巢原理]
把 $n+1$ 只鸽子放进 $n$ 个鸽巢，至少有一个巢里有两只或更多鸽子。
\end{dingli}

\begin{proof}
反设每个巢里至多一只鸽子，那么 $n$ 个巢至多装 $n$ 只，与共有 $n+1$ 只矛盾。
\end{proof}

一句话的定理，能解决看上去完全无从下手的问题。例如：全校 400 名学生中，必有两人同一天生日——因为生日最多 366 种，学生比“巢”多。再如：任意 6 个人中，必有 3 人互相认识，或者 3 人互不认识。这个结论可以用鸽巢原理证：任取一人甲，甲与其余 5 人的关系只有“认识、不认识”两种，5 只鸽子进 2 个巢，必有一种关系至少涉及 3 人，再继续讨论这 3 人之间的关系即可（细节留给思考题）。鸽巢原理的深层含义是：\emph{当对象的个数超过分类的格数时，碰撞不可避免}。这个朴素思想是现代密码学中“哈希碰撞”研究和计算机科学中“抽屉论证”的原型。

\section{把每支球队画成一个点：图、树与路}

现在换一个角度看单循环赛。20 支球队两两比赛，我们关心的其实不是“谁赢谁输”，而是“谁和谁有关系”。把每支球队画成一个点，每场比赛在两点间连一条线，整个赛程就变成一张“点线图”。生活中的同类结构到处都是：城市和航线、人和好友关系、网页和超链接、路口和街道、芯片上的元件和导线。数学家把这类结构抽象成一个概念。

\begin{dingyi}[图]
一个\emph{图}由一组\emph{顶点}和一组\emph{边}组成，每条边连接两个顶点。如果边有方向（比如单行道），称为有向图；如果边上标有数字（比如路程、费用），称为加权图。与某个顶点相连的边的条数，叫作该顶点的\emph{度数}。
\end{dingyi}

注意这个定义故意扔掉了所有几何信息：图里的“点”没有位置，“线”没有长短和曲直，唯一的信息是\emph{谁和谁连着}。这是一种故意的遗忘——忘掉不重要的，才能看清结构本身。

\begin{center}
\begin{tikzpicture}[every node/.style={circle,draw,inner sep=1.5pt,minimum size=6mm}]
\node (a) at (0,2) {甲};
\node (b) at (3,2.5) {乙};
\node (c) at (5,1) {丙};
\node (d) at (3,-0.5) {丁};
\node (e) at (0.5,-0.5) {戊};
\draw (a)--(b) (b)--(c) (c)--(d) (d)--(e) (e)--(a) (a)--(c);
\end{tikzpicture}
\end{center}

上面这张小图里，顶点甲的度数是 3（连着乙、丙、戊），顶点戊的度数是 2。有一个所有图都遵守的定律。

\begin{dingli}[握手引理]
在任何图中，所有顶点的度数之和等于边数的两倍。
\end{dingli}

\begin{proof}
逐条边去数。每条边有两个端点，每数一条边，就给它的两个端点的度数各贡献 1。所以度数总和恰好是边数的 2 倍。
\end{proof}

这个定理有一个立刻可用的推论：任何图中，\emph{度数为奇数的顶点必有偶数个}。因为度数总和是偶数，奇数加奇数才可能是偶数，奇度顶点必须成对出现。翻译回联欢会：任何一场聚会中，“握手次数为奇数”的人数一定是偶数。读者下次参加聚会可以默默验证——这不是巧合，是定理。

图论里最干净的一类图是\emph{树}：连通（任意两点之间都有路可走）且没有圈的图。树像真实的大树一样，从任何一个点出发都不会绕回原地。一个基本事实是：$n$ 个顶点的树恰好有 $n-1$ 条边——每条边都“必不可少”，删去任何一条，图就断成两截。家族谱系图、电脑里的文件夹目录、淘汰赛的晋级图，都是树。树的一个重要用途是“用尽量少的边把大家连通”：要在 6 个村庄之间修路，让所有村庄互通，最少修 5 条；如果想让总造价最低，就要在一堆候选边里挑出一棵“最便宜的树”，这就是最小生成树问题，它有简单高效的贪心算法。

更有现实气息的问题是\emph{最短路}。给一张加权图（顶点是路口，边是街道，数字是长度），从家到学校走哪条路总长度最小？读者可能想说“枚举所有路线取最短”，但路口一多，路线条数会像阶乘一样爆炸。1950 年代，计算机科学家戴克斯特拉（Dijkstra）给出了一个贪心算法：从起点出发，始终优先处理“当前已知最近”的那个顶点，用它去刷新邻居们的距离记录，像水波一样一圈圈向外扩散，直到抵达终点。这个算法的每一步只做局部比较，却保证找到全局最短路——前提是所有边的长度非负。今天导航软件里的路径规划，核心仍然是这个思想的加强版。

把“长度”换成“容量”，就进入另一个世界：\emph{网络流}。把图想象成水管网，每条边有最大流量（口径），问：从水厂到你家，单位时间最多能送多少水？直观上，答案受“瓶颈”限制：如果在地图上画一条线把水厂和你家分开，所有跨过这条线的水管的总容量，就是这条线能“放行”的上限。所有这样的分割线中，最紧的那条决定了真实上限。这就是著名的\emph{最大流最小割定理}：网络中的最大流量，等于最小割（容量最小的分割线）的容量。我们不证明它，但请读者品味它的形式：一个“最大化”问题的答案，恰好等于另一个“最小化”问题的答案。这种“最大等于最小”的对偶现象，在整个数学里反复出现，每出现一次都意味着找到了结构深处的对称。快递分拣、电网调度、甚至给球队安排赛程，背后都有网络流的影子。

\section{算法：让机器按图索骥}

到目前为止，问题都是我们用手和脑解决的。可现实里的图有上百万个顶点——全国的路网、整个互联网的链接结构。这时候需要把解题方法写成\emph{算法}：一份严格、机械、任何人（或任何机器）照着做都能得到答案的步骤清单。算法是离散数学与计算机科学的接头处。

设计算法有两种基本姿势，本章前面已经偷偷见过。第一种是\emph{递归}：把问题拆成同类型的更小问题，直到小到直接有答案。爬楼梯走法 $f_n=f_{n-1}+f_{n-2}$ 就是递归；电脑里“遍历文件夹”也是递归——“列出这个文件夹的内容”就是“对每个子文件夹再执行列出”。第二种是\emph{搜索}：把所有可能性组织成一棵“决策树”，然后系统地走这棵树。走迷宫就是搜索：每个岔路口是一个顶点，每条走廊是一条边，“找到出口”就是在这张图上找一条路。

搜索又分两种性格。\emph{深度优先}像一根筋的探险者：沿一条路走到底，撞墙再回头；\emph{广度优先}像一圈圈扩散的水波：先检查所有一步可达的位置，再检查两步可达的，最短路问题里它保证先找到的出口一定是最近的。戴克斯特拉算法本质上就是广度优先在加权图上的升级版。

但是，“能算出答案”和“算得起答案”是两回事。这就需要本章最后一件武器：\emph{算法复杂度}的直观说法。

\begin{shuyu}{大 O 记号}
\textbf{直观解释}：只关心当问题规模 $n$ 变大时，算法的工作量大约按什么方式增长，忽略常数倍和细枝末节。

\textbf{正式说法}：若工作量不超过 $n$ 的某个固定倍数，记作 $O(n)$；不超过 $n^2$ 的固定倍数，记作 $O(n^2)$；依此类推。

\textbf{一个例子}：在 $n$ 个数里找最大值的算法是 $O(n)$——每个数看一眼就够；枚举 $n$ 个人的所有排队方案是 $O(n!)$——根本跑不完。
\end{shuyu}

用大 O 的眼光重看本章：在 $n$ 个顶点的图上跑戴克斯特拉算法，工作量大约是 $O(n^2)$，路口数扩大 10 倍，时间大约扩大 100 倍——可以接受。检查“$n$ 个人两两关系”的问题是 $O(n^2)$——也可以接受。但如果一个算法的复杂度是 $O(2^n)$ 或 $O(n!)$，故事就完全不同了：$2^{40}$ 已经超过一万亿，$25!$ 更是天文数字。规模只增长到几十，机器就要跑到天荒地老。这类问题里有不少著名难题，比如“旅行商问题”：给若干城市，找一条走遍所有城市再回到起点的最短回路。至今没有人知道有没有本质上比“近乎枚举”快得多的通用算法——这是计算机科学最大的未解之谜之一（通常写作“P 与 NP 问题”），数学家们悬赏百万美元等待答案。

于是离散数学给我们上的最后一课是双面的：一方面，图、递推和贪心让我们能把看似不可能的大问题压进可计算的范围；另一方面，复杂度理论冷酷地标出了边界——有些问题不是“暂时算不动”，而可能是“本质上就算不动”。知道哪边是哪边，本身就是深刻的数学知识。

\section{离散结构、密码与数论：为什么它们分不开}

本章的主题看似是“数数和连线”，但它和第 13 章的数论、和你手机里的密码，是同一条战壕里的战友。理由有三层。

第一层，\emph{计算机本身就是离散机器}。它存的是 0 和 1 的串，走的是一步一条指令，天然生活在离散世界里。连续的量（比如圆的面积）进了计算机也要被切碎成有限位小数。所以研究离散结构的数学——组合、图、逻辑、递推——就是计算机科学的“母语教材”。

第二层，\emph{密码的安全性来自离散难题}。第 13 章讲过，把两个大质数相乘很容易，把乘积拆回两个质数极难——这是 RSA 密码的地基。请注意这个不对称正是“复杂度”问题：正向是 $O(\text{多项式})$ 的乘法，逆向目前只知道接近枚举的慢算法。离散对数、格上的最短向量等离散结构难题，也都站在现代密码的底层。本章的鸽巢原理则在另一头守护密码学：哈希函数把任意长的消息压成固定长度的“指纹”，消息无穷而指纹有限，鸽巢原理保证碰撞必然存在——密码学家要做的，是让\emph{找到}碰撞这件事在计算上不可行。

第三层，\emph{数论本身也是离散数学}。整数是一个一个的，质数是一颗一颗散落在数轴上的点。组合工具一直在为数论服务：二项式系数里藏着质数的行踪（比如当 $p$ 是质数时，$\binom{p}{1},\binom{p}{2},\dots,\binom{p}{p-1}$ 全都被 $p$ 整除，读者可以用 $p=5$ 亲手验证）；递推数列的整除性质是竞赛和研究的常客；图论甚至被用来研究质数之间的“距离”。

\begin{lianjie}
\textbf{向前看}：本章的乘法原理和阶乘，会在第 18 章的概率里再次登场——“等可能的样本空间”其实就是一次组合计数；图的连通性和“圈”的概念，会在第 16 章拓扑学里被推广成“欧拉示性数”；算法与不可计算的边界，会在第 19 章继续深挖。

\textbf{向后看}：第 7 章的矩阵可以记录一张图的连接关系（邻接矩阵），矩阵乘法一次，数出的就是“两步可达”的路径条数——线性代数和图论在这里握手。第 13 章的同余和质数，是本章组合数整除性质和密码应用的共同地基。
\end{lianjie}

\section{数学实验}

\begin{shiyan}
\textbf{实验一：数树。} 在纸上画 4 个顶点，试着画出所有“不同构”（连法本质上不一样）的树，数一数有几种；再画 5 个顶点数一数。（4 个顶点有 2 种，5 个顶点有 3 种。）体会一下：这个数列增长得很快，到 10 个顶点时已有 106 种。

\textbf{实验二：手算最短路。} 画一张 5 个路口的小图：家、学校、公园、书店、车站，随意连七八条边并标上长度。用戴克斯特拉的思想，从“家”开始，把每个路口旁标注“目前已知的最近距离”，每处理一个路口就刷新它的邻居。算完和自己凭眼睛猜的路线比一比。

\textbf{实验三：二项式系数与质数。} 算出 $\binom{5}{1},\binom{5}{2},\binom{5}{3},\binom{5}{4}$，检查是否都被 5 整除；再对 $n=6$ 做同样的事，看看结论还在不在。想一想差别在哪。
\end{shiyan}

\begin{toolbox}
本章收入工具箱的新工具：
\begin{itemize}
\item \textbf{加法原理与乘法原理}：分类用加，分步用乘；先问“我数重了吗”。
\item \textbf{排列与组合}：有序用 $A_n^k=\dfrac{n!}{(n-k)!}$，无序用 $\dbinom{n}{k}=\dfrac{n!}{k!(n-k)!}$。
\item \textbf{递推}：用更小规模的同类问题定义当前问题，如 $f_n=f_{n-1}+f_{n-2}$。
\item \textbf{二项式定理}：$(a+b)^n$ 的展开系数就是组合数；令 $a=b=1$ 得子集总数 $2^n$。
\item \textbf{鸽巢原理}：鸽子比巢多，必有碰撞——碰撞不可避免性论证的原型。
\item \textbf{图的语言}：顶点、边、度数；握手引理 $\sum\text{度数}=2\times\text{边数}$。
\item \textbf{树、最短路、网络流}：连通且省边的骨架、贪心扩散、最大流等于最小割。
\item \textbf{大 O 记号}：用增长方式（$O(n)$、$O(n^2)$、$O(2^n)$、$O(n!)$）判断算法跑不跑得动。
\end{itemize}
\end{toolbox}

\begin{pandeng}
继续向前可以攀登的方向：
\begin{itemize}
\item \textbf{生成函数}：把数列 $a_0,a_1,a_2,\dots$ 压缩成一个函数，用代数方法解递推——斐波那契数列的通项公式就是这样求出来的。
\item \textbf{拉姆齐理论}：“6 人中必有 3 人互识或互不认识”只是开端；一般地，足够大的结构里必然出现足够整齐的子结构，这是“完全的无序不可能”的数学。
\item \textbf{P 与 NP}：可验证与可求解之间的鸿沟，克雷数学研究所的七大千禧年难题之一。
\item \textbf{谱图论}：用矩阵的特征值研究图的连通性，支撑了搜索引擎的网页排序。
\item \textbf{极值图论}：在“不能有三角形”之类的禁令下，一张图最多能有多少条边。
\end{itemize}
\end{pandeng}

\section*{思考题}

\textbf{观察题}
\begin{sikao}
\item 计算 $\binom{10}{2}$ 和 $\binom{10}{8}$，你会发现它们相等。用“选 2 个人去开会”和“选 8 个人留下”的关系解释为什么 $\binom{n}{k}=\binom{n}{n-k}$ 永远成立。\emph{提示：每一次“选出 $k$ 个”都同时“留下 $n-k$ 个”。}
\item 数一数你所在小组（或家里）每个人认识组内其他人的数目，把度数加起来，看看是不是偶数，并用握手引理解释。\emph{提示：把“认识关系”画成图。}
\end{sikao}

\textbf{证明题}
\begin{sikao}
\item 证明：$n$ 个顶点的树恰好有 $n-1$ 条边。\emph{提示：对顶点数用归纳法；一棵树总能找到一片“叶子”（度数为 1 的顶点），摘掉它和连着的边，剩下的仍是一棵更小的树。}
\item 证明：任意 6 个人中，必有 3 人两两互相认识，或 3 人两两互不认识。\emph{提示：正文已经走了一半——固定一人，用鸽巢原理找出与他关系相同的 3 人；再分“这 3 人中有两人关系也如此”和“没有”两种情况。}
\end{sikao}

\textbf{挑战题}
\begin{sikao}
\item 在一条直线上等距取 $n$ 个点，把每对点连一条线段。所有线段的中点中，至少有多少个互不相同的点？\emph{提示：先对 $n=3,4,5$ 画图数一数；把点放在坐标 $1,2,\dots,n$ 上，中点坐标是 $\dfrac{i+j}{2}$，考察它能取哪些值。}
\item 8 枚外观相同的金币中有 1 枚较轻的假币，用一架没有砝码的天平，至少称几次能保证找出假币？如果改成“不知道假币偏轻还是偏重”，12 枚金币又需要几次？\emph{提示：每称一次有三种结果（左沉、右沉、平衡），$k$ 次称量最多区分 $3^k$ 种情形——把每一次称量看作决策树上的一个分叉。}
\end{sikao}

\section*{通往下一章的门}

这一章我们学会了数离散的物件：排队的人、握手的次数、图里的边。我们一再看到同一个现象：离散对象的局部规律（比如 $\binom{p}{k}$ 被质数 $p$ 整除）往往由质数在幕后操纵。可是，如果我们把镜头拉远，问一个宏观问题——质数这些“离散的点”，在整条数轴上到底铺得多密？

具体地说：不超过 100 的质数有 25 个，不超过 1000 的有 168 个，不超过一百万的有 78498 个。这串数字有没有公式？越往后质数越稀疏，但稀疏的速度是匀速的吗？有趣的是，回答这个纯离散的问题，最强有力的工具竟然来自最连续的世界——微积分、无穷级数，甚至复数平面上的函数。数学家们发现，质数的分布仿佛被一条“连续的影子”精确牵引着，而这条影子的秘密，浓缩在一个叫黎曼 $\zeta$ 函数的对象身上。它身上还挂着数学界最著名的未解之谜。

离散与连续，这两条看起来背道而驰的河流，将在下一章交汇。
